將大型語言模型（LLM）適配到新語言是一個成本高且不透明的過程。理解語言模式如何習得新語言和多語言能力是實現高效率適配的關鍵。以往關於多語言可解釋性的研究主要集中在訓練後的模型如何處理多語言指令，而忽略了它們在訓練過程中習得新語言的機制。我們以僅包含解碼器的Transformer模型為例，從兩個功能性認知特化的角度－語言感知（輸入理解）和語言生成（輸出產生）－來研究這些訓練動態。透過在低資源語言上進行實驗，我們展示了感知和語言生成特化如何在語言模型的不同區域中湧現，方法是沿著模型的輸入和輸出方向進行層消融掃描。基於觀察到的特化模式，我們提出了CogSym，一種逐層啟發式方法，它透過僅微調少數幾個早期和晚期層來實現有效的適配。我們證明，僅調整最外層的25\%即可使下游任務的性能與完全微調的基線性能偏差控制在2-3\%以內。此外，我們也從第一語言習得（FLA）和第二語言習得（SLA）的角度分析了語言學習模型（LLM）中的語言學習困難。我們的研究揭示了目前語言學習模型架構與人類語言習得過程中的認知和行為之間的相似之處。

\chinesekeywords{大型語言模式, 多語現象, 語言習得}